\begin{abstract}

Multilingual meetings require speech-to-speech translation that preserves each speaker's contributions, even during overlapping speech. In practice, translated speech can queue behind earlier utterances, delaying what listeners hear after the text has been committed. However, existing systems primarily handle one speaker at a time and focus on when translated text is produced rather than when its speech is heard. To address these limitations, we introduce Meeting Interpretation with Speaker Tracking (\sys{}), a speaker-attributed simultaneous translation framework with playback-aware speech generation. \sys{} interleaves acoustic tokens with speaker-activity representations from a frozen simultaneous diarizer, enabling separate, time-stamped translations of overlapping speakers' contributions. To reduce listener-side delay, we introduce playback alignment, a linear-chain conditional random field (CRF) that scores candidate speech renderings by their likelihood, recognition fidelity, and the delay they impose on later speech. Exact forward--backward inference and posterior regularization yield a latency-constrained distribution that is distilled into the speech generator, allowing translated segments to be spoken in each speaker's voice. We introduce cpASR-BLEU and tcpASR-BLEU to evaluate speaker-attributed synthesized speech from ASR transcripts, with and without temporal constraints. Across AMI and AliMeeting, \sys{} improves tcpASR-BLEU by 1.8–12.6\% and reduces speech completion delay by 16.4–29.3\% over a baseline with identical text. To our knowledge, MIST is the first speaker-attributed simultaneous speech-to-speech translation system to translate multi-talker meetings, including overlapping speech, in each speaker's voice, and the first to evaluate speaker-attributed speech quality with novel metrics and latency over complete meetings.


\end{abstract}
